\chapter*{Interlude: A Complete Verification, Entirely by Hand}
\markboth{Interlude: A Complete Verification, Entirely by Hand}{}
\addcontentsline{toc}{chapter}{Interlude: A Complete Verification, Entirely by Hand}

Before the full field campaign of the next chapter, we pause to do
something no other book chapter will ask of you again: verify an
arithmetic implementation \emph{completely}, on paper, with nothing left
to the machine --- and then watch the same proof re-enacted in Lean,
line by line. The system is the two-limb miniature from
Chapter~\ref{ch:denotation}'s exercises; the point is that after this
interlude, nothing in the real proofs is a new \emph{kind} of thing ---
only a bigger instance of what your own pen has already done.

\section*{I.1 \quad The system, restated in full}

\begin{itemize}[leftmargin=1.4em]
\item \textbf{Representation.} A value is a pair of natural-number limbs
$(a_0, a_1)$, intended as base-4 digits.
\item \textbf{Bounds invariant.} $\mathrm{Bnd}(a) \;:\Longleftrightarrow\;
a_0 < 4 \,\wedge\, a_1 < 4$.
\item \textbf{Denotation.} $\denote{a} = a_0 + 4 a_1 \in \Zmod{15}$.
The modulus $15$ is chosen so that the radix squared folds:
$16 \equiv 1 \pmod{15}$ --- the toy twin of $2^{255} \equiv 19$.
\item \textbf{The implementation under verification.}
\[
\code{add}(a, b) := \big(\, s_0 \bmod 4 + \lfloor s_1 / 4 \rfloor,\;\;
s_1 \bmod 4 \,\big),
\quad\text{where } s_0 = a_0 + b_0,\;
s_1 = a_1 + b_1 + \lfloor s_0 / 4 \rfloor .
\]
In words: add low limbs; carry into the high sum; the high sum's own
overflow --- weight $16 \equiv 1$ --- folds back into the low limb.
\end{itemize}

\section*{I.2 \quad The specification}

The two-clause shape of Chapter~\ref{ch:denotation}, in full:
\[
\textbf{(Spec)}\qquad
\mathrm{Bnd}(a) \to \mathrm{Bnd}(b) \to
\underbrace{\big( \code{add}(a,b)_0 < 5 \,\wedge\, \code{add}(a,b)_1 < 4 \big)}_{\text{bounds clause}}
\;\wedge\;
\underbrace{\denote{\code{add}(a,b)} = \denote{a} + \denote{b}}_{\text{value clause}} .
\]
Pause on the asymmetric output bound: the low limb is only promised
$< 5$, not $< 4$ --- the fold can push it one past a clean digit. An
honest spec records where the envelope actually lands, and the next
operation's hypotheses must accept what this one delivers. (In the real
field code this is the $2^{51}$-in, $2^{52}$-out pattern of
Chapter~\ref{ch:rust}'s worked example.)

Run the Chapter~\ref{ch:honesty} substitution test before proving
anything: would ``return $(0,0)$ always'' pass? Bounds clause: yes.
Value clause: no --- $a = (1,0), b = (0,0)$ gives $1 \neq 0$. The spec
has teeth. Now we earn it.

\section*{I.3 \quad The bounds clause, proved by hand}

Assume $\mathrm{Bnd}(a)$, $\mathrm{Bnd}(b)$; so
$a_0, a_1, b_0, b_1 \le 3$. Chase every intermediate through its
interval:
\[
\begin{array}{lclcl}
s_0 &=& a_0 + b_0 &\le& 3 + 3 = 6\\[2pt]
\lfloor s_0/4 \rfloor &\le& \lfloor 6/4 \rfloor &=& 1\\[2pt]
s_1 &=& a_1 + b_1 + \lfloor s_0/4 \rfloor &\le& 3 + 3 + 1 = 7\\[2pt]
\lfloor s_1/4 \rfloor &\le& \lfloor 7/4 \rfloor &=& 1\\[2pt]
{\code{add}}(a,b)_0 &=& s_0 \bmod 4 + \lfloor s_1/4 \rfloor &\le& 3 + 1 = 4 \;<\; 5 \quad✓\\[2pt]
{\code{add}}(a,b)_1 &=& s_1 \bmod 4 &\le& 3 \;<\; 4 \quad✓
\end{array}
\]
Done --- and notice the proof \emph{is} the headroom audit of
Chapter~\ref{ch:why}, in miniature: every line asks ``how big can this
intermediate get,'' and the two ✓ lines are the treaty the spec
promised. Notice also that the bound $<5$ is \emph{tight}: at
$a = b = (2, 3)$, $s_0 = 4$ folds a carry and $s_1 = 7$ folds another,
giving low limb $0 + 1 = 1$\dots\ try $a = b = (3,3)$: $s_0 = 6$,
$s_1 = 7$, output $(2 + 1, 3) = (3,3)$ --- hmm, low limb $3$. Exercise
I.1 asks you to find the input that actually achieves $4$; it exists,
and finding it will teach you more about tightness than ten theorems.

\section*{I.4 \quad The value clause, proved by hand}

The pathway (identical to the one you will use in Lean): first prove an
\emph{exact integer identity} --- no mod, no hand-waving --- then let the
clock face absorb the multiple of $15$.

\textbf{Step 1: the division algorithm, twice.} For any naturals,
$x = 4\lfloor x/4 \rfloor + (x \bmod 4)$. Apply to $s_0$ and $s_1$:
\[
s_0 = 4\lfloor s_0/4 \rfloor + (s_0 \bmod 4),
\qquad
s_1 = 4\lfloor s_1/4 \rfloor + (s_1 \bmod 4).
\]

\textbf{Step 2: compute the output's integer value, plus a correction.}
Write $c_1 := \lfloor s_0/4 \rfloor$ (low carry) and
$c_2 := \lfloor s_1/4 \rfloor$ (the folded top carry). The output's
positional value is
\[
\begin{array}{lcl}
\denote{\code{add}(a,b)}_{\Z}
&=& \big( s_0 \bmod 4 + c_2 \big) + 4\,(s_1 \bmod 4)\\[3pt]
&=& \big( s_0 - 4c_1 + c_2 \big) + 4\,( s_1 - 4 c_2 )
    \qquad\text{(Step 1, both equations)}\\[3pt]
&=& s_0 - 4c_1 + 4 s_1 - 15\, c_2\\[3pt]
&=& s_0 - 4c_1 + 4\,( a_1 + b_1 + c_1 ) - 15\,c_2
    \qquad\text{(definition of } s_1\text{)}\\[3pt]
&=& (a_0 + b_0) + 4\,(a_1 + b_1) \;-\; 15\, c_2 .
\end{array}
\]
Watch what happened in the last two lines: the low carry $c_1$
\emph{cancelled exactly} --- $-4c_1$ against $+4c_1$ --- because a carry
is value moved between positions, not value created; and the top carry
survived only as $-15\,c_2$, because folding weight $16$ down to weight
$1$ loses exactly $15$ per unit. The bookkeeping \emph{is} the
arithmetic meaning of ``carry'' and ``fold,'' laid bare.

\textbf{Step 3: pass to the clock face.} Reducing mod $15$, the term
$15\,c_2$ vanishes:
\[
\denote{\code{add}(a,b)}
\;=\; (a_0 + 4a_1) + (b_0 + 4b_1)
\;=\; \denote{a} + \denote{b} \quad\text{in } \Zmod{15}. \qquad\blacksquare
\]
The value clause is proved --- for \emph{every} bounded input, by three
steps of eighth-grade algebra. You have now verified an arithmetic
implementation by hand, completely: no case was sampled, no input
untested, because no input was \emph{tested} at all. The quantifier came
from algebra, not enumeration.

\section*{I.4b \quad Multiplication's fold, by the same method}

One more clause, to see the method survive a second operation. The toy
multiplication (as in \code{exercises/Ch09.lean}) computes the three
schoolbook columns and folds the top one:
\[
\code{mulVal}(a,b) := \underbrace{a_0 b_0}_{c_0} +
\underbrace{a_1 b_1}_{c_2\ \text{(folded)}} +
4\,\underbrace{(a_0 b_1 + a_1 b_0)}_{c_1},
\]
claim: $\denote{\code{mulVal}(a,b)} = \denote{a}\cdot\denote{b}$ ---
this time with \emph{no bounds hypotheses at all}. Same three-step
pathway, but Step 2's identity is now pure algebra. Expand the
right-hand side:
\[
(a_0 + 4a_1)(b_0 + 4b_1)
= a_0 b_0 + 4\,(a_0 b_1 + a_1 b_0) + 16\, a_1 b_1 .
\]
Compare with $\code{mulVal}$: the only mismatch is the top column's
weight --- $16\,a_1 b_1$ there, $1 \cdot a_1 b_1$ here. So the exact
integer identity is
\[
\code{mulVal}(a,b) + 15\, a_1 b_1 \;=\; (a_0 + 4a_1)(b_0 + 4b_1),
\]
--- verify it by cancelling the displayed expansions --- and Step 3
kills the $15\,a_1 b_1$ on the clock face. $\blacksquare$

Two observations before we translate. First, \emph{why} no bounds were
needed: this theorem speaks only of values, and positional value is
indifferent to digit discipline; bounds enter only when machine words
must contain the intermediates (the toy \code{mulVal} returns a bare
natural number --- the real \code{mul} returns limbs, and \emph{re-limbing
is where its bounds clauses live}). Second, the correction term's shape:
$15$ per unit of \emph{folded column}, exactly the ``cost of a fold''
law from I.4 --- one law, two operations, and in the real system the
same law reads ``$p$ per folded column,'' four times over
(Chapter~\ref{ch:denotation}'s fold worked example).

\section*{I.5 \quad The same proof, in Lean, line by line}

Here is the machine version (it is
\code{solutions/Ch09.lean}, compiled and axiom-audited), interleaved
with the paper steps it enacts:

\begin{center}
\small
\begin{tabular}{@{}p{0.47\linewidth}p{0.47\linewidth}@{}}
\toprule
\textbf{Paper (sections I.3--I.4)} & \textbf{Lean} \\
\midrule
unpack the four digit bounds &
\lean{obtain ⟨ha1, ha2⟩ := ha} \dots \\
\addlinespace[3pt]
the interval chase of I.3 &
\lean{simp only [add]; omega} \\
\addlinespace[3pt]
Steps 1--2: the exact integer identity, carries cancelling ---
stated with the $15 c_2$ correction on the left &
\lean{have key : (add a b).1 + 4*(add a b).2 + 15*(...) = ...  := by simp only [add]; omega} \\
\addlinespace[3pt]
Step 3: cast to $\Zmod{15}$; the $15$ becomes $0$ &
\lean{push_cast at this;} \lean{rw [show (15 : ZMod 15) = 0 by decide]} \\
\addlinespace[3pt]
``\dots{} $= \denote{a} + \denote{b}$. $\blacksquare$'' &
\lean{simp only [denote]; linear\_combination this} \\
\bottomrule
\end{tabular}
\end{center}

Read the correspondence twice, because it is the book's whole method in
one table. The interval chase you did line-by-line is one \lean{omega}
call --- the machine \emph{decides} what you \emph{derived}, which is why
Chapter~\ref{ch:automation} called these procedures contracts. The
delicate part on paper --- Step 2's cancellation --- is delicate in Lean
too: it is the \lean{key} identity, and getting its statement right
(which correction term? on which side?) is where the human does the
thinking in both media. Step 3 is bookkeeping in both media. The
division of labor between you and the machine is the same division
between insight and verification your paper proof already had.

And now the punchline of the whole interlude: the field-layer proof of
\code{dalek-ed25519-verified} is \emph{this proof} with five limbs
instead of two, radix $2^{51}$ instead of $4$, fold constant $19$
instead of $1$ (times the fold identity's bookkeeping), and five carry
cancellations instead of one. Every step type --- interval chase, exact
identity with correction term, cast and kill the modulus --- you have
now executed by hand.

\section*{Interlude exercises}

\noindent{\bfseries\color{accent}Exercise I.1.}\ Find bounded inputs
achieving the tight output bound: $\code{add}(a,b)_0 = 4$. (Systematic
pathway: you need $s_0 \bmod 4 = 3$ \emph{and} a fold, $c_2 = 1$; work
out what $s_0$ and $s_1$ must be, then pick digits realizing them.)

\smallskip
\noindent{\bfseries\color{accent}Exercise I.2.}\ Re-run the \emph{entire}
I.3--I.4 proof for the variant system with radix $8$ and modulus $63$
(so $64 \equiv 1$): two limbs $< 8$, denotation $a_0 + 8a_1$. Every step
should transpose mechanically --- when one doesn't, you have found a
place where you were pattern-matching instead of understanding; welcome,
that is the exercise working.

\smallskip
\noindent{\bfseries\color{accent}Exercise I.3.}\ (Bridge to the real
thing) In the I.4 cancellation, the low carry vanished and the top carry
cost $15$ per unit. For the real radix-51 system: state what the
analogous ``cost per unit of top carry'' is, and verify your answer
against the fold identity of Chapter~\ref{ch:denotation}'s worked
example.

\smallskip
\noindent{\bfseries\color{accent}Exercise I.4.}\ (The missing clause)
Section I.4b's \code{mulVal} returns a bare number, not limbs --- so its
theorem has a value clause but no bounds clause. Complete the design on
paper: define $\code{relimb}(v) := (v \bmod 4 + \text{(fold)},\, \dots)$
--- that is, re-express a value $v \le 90$ (the largest \code{mulVal}
of bounded inputs --- verify this bound first!) as two limbs using the
division algorithm and one fold. How many fold rounds does $v \le 90$
need before both limbs are $< 4$? State the two-clause spec your
$\code{relimb}$ would carry.

\section*{Interlude solutions and pathways}
\solutionsintro

\solhead{I.1}
\pathway Work backwards from the two requirements. $c_2 = 1$ needs
$s_1 \ge 4$; $s_0 \bmod 4 = 3$ with a useful range of $s_0$: since
$s_0 \le 6$, ``$s_0 \bmod 4 = 3$'' means $s_0 = 3$ (no carry out) --- so
the fold must come entirely from $a_1 + b_1$.
\answer Take $s_0 = 3$ (say $a_0 = 3, b_0 = 0$), no low carry; then
$s_1 = a_1 + b_1 \ge 4$ (say $a_1 = b_1 = 2$, $s_1 = 4$): output low
limb $= 3 + 1 = 4$. Concretely $a = (3,2)$, $b = (0,2)$:
$\code{add} = (4, 0)$. Check the value clause holds anyway:
$\denote{(4,0)} = 4$ and $\denote{a} + \denote{b} = 11 + 8 = 19 \equiv 4$
✓ --- the bound is tight but the meaning is intact, which is exactly why
the spec's bounds clause said $< 5$ and not $< 4$. (If you tried to
strengthen the spec to $< 4$, this input is the counterexample the
failed proof would point at --- Chapter~\ref{ch:pyramid}'s graduation
exercise mechanism, met early.)

\solhead{I.2}
\pathway Transcribe I.3--I.4 replacing $(4, 15, 16)$ by $(8, 63, 64)$
and the digit bound $3$ by $7$.
\answer Bounds chase: $s_0 \le 14$, $c_1 \le 1$, $s_1 \le 15$,
$c_2 \le 1$, output low $\le 7 + 1 = 8 < 9$, output high $\le 7 < 8$.
Value identity: output value
$= s_0 - 8c_1 + 8s_1 - 63\,c_2
= (a_0 + b_0) + 8(a_1 + b_1) - 63\,c_2$ --- low carry cancels
($-8c_1 + 8c_1$), top carry costs $64 - 1 = 63$ per unit, which
vanishes mod $63$. Every step transposed; the only places requiring
thought were the two spots where the old numbers were \emph{related}
($15 = 16 - 1$ became $63 = 64 - 1$; $\lfloor 14/8 \rfloor = 1$ needed
re-checking, not copying) --- and those two spots are precisely the
design constraints of the system: fold constant $=$ radix$^2 - $modulus
relation, and single-bit carries. Now you know which numbers in such a
system are load-bearing.

\solhead{I.3}
\pathway In I.4 the cost was (weight folded from) $-$ (weight folded
to) $= 16 - 1 = 15$ per unit of top carry. Transpose the weights.
\answer The real system folds weight $2^{255}$ down to weight $2^{0}$
\emph{with a factor of 19}: one unit of top overflow re-enters as $+19$
instead of $+2^{255}$, so the correction is
$2^{255} - 19 = p$ per unit --- the cost is exactly one prime $p$, which
is why it vanishes mod $p$, and why the fold constant \emph{must} be
$19$ and nothing else: it is the unique value making the correction a
multiple of the modulus. Cross-check with the fold worked example:
$2^{51(k+5)} \equiv 19 \cdot 2^{51k}$ says each folded column's
correction is $2^{51k}(2^{255} - 19) = 2^{51k}\, p$ ✓. The toy's
``$15$ per carry'' and the field's ``$p$ per carry'' are the same
sentence at different type sizes --- carry the sentence with you into
the next chapter.

\solhead{I.4}
\pathway Bound \code{mulVal} first by maximizing each column, then
design \code{relimb} as ``split by the division algorithm, fold the
overflow, repeat until in range,'' tracking the bound through each
round exactly as I.3's chase did.
\answer Bound: columns give $\le 9 + 9 + 4 \cdot 18 = 90$ ✓. One
\code{relimb} round on $v$: limbs $(v \bmod 4 + \lfloor v/16 \rfloor,\;
\lfloor v/4 \rfloor \bmod 4)$ --- value preserved mod $15$ since the
fold costs $15$ per unit (I.4's law). Chase the bound: round one sends
$v \le 90$ to $v' \le 3 + 5 + 12 = 20$; a value $< 16$ needs no fold
and splits into two clean digits, so only $v' \in \{16, \dots, 20\}$
needs a further round, which lands $\le 3 + 1 + 4 = 8 < 16$ ---
\textbf{at most three rounds}, usually two. The spec your \code{relimb}
carries is precisely the two-clause shape: (bounds) both output limbs
$< 4$; (value) $\denote{\code{relimb}(v)} = v$ in $\Zmod{15}$. And now
compose: $\code{relimb} \circ \code{mulVal}$ has a full multiplication
spec, assembled from two lemmas --- which is, in miniature, exactly how
the real \code{mul} proof is architected (fold theorem + carry-chain
re-limbing), and why Chapter~\ref{ch:field}'s campaign put
\code{reduce} before \code{mul}.

\begin{checkpoint}
You have now: proved both clauses of a two-clause spec with your own
pen; watched each paper step map onto one Lean tactic; transposed the
proof to a sibling system and located its load-bearing constants; and
priced a fold in multiples of the modulus. The next chapter's campaign
is this interlude at production scale --- when it feels big, return
here and find the step type you are on.
\end{checkpoint}
